\section{Fixed-column product-ready reassembly}
\label{sec:tpc36-fixed-column}

The global theorem separates the two row labels, but a global factor may
still depend jointly on the source \(\ell\) and divisor coordinate \(d\)
through \(m=\ell d\).  The product--slope operator needs one further
separation.  Orbit slicing provides the correct order: fix the opposite
row first.

\subsection{The exact Ferrers formula}

Fix
\begin{equation}
 \gamma=(\ell_\gamma,e),\qquad n=m_\gamma=\ell_\gamma e.
 \label{eq:tpc36-fixed-column}
\end{equation}
Let \(\cD=\{d_1<\cdots<d_{N_D}\}\) be an ambient integer set containing
the physical \(d\)-support.  One may take the full dyadic integer interval,
so
\begin{equation}
 N_D\ll D.
 \label{eq:tpc36-divisor-count}
\end{equation}
Squarefreeness, coprimality, and smooth support are retained later as
unary \(d\)-weights.  Write \(r(d_k)=k\).  For every source \(\ell\), put
\begin{align}
 k_-^\gamma(\ell)
 &:=\#\{d\in\cD:\ell d<n-\Delta\},
 \label{eq:tpc36-column-lower-rank}\\
 k_+^\gamma(\ell)
 &:=\#\{d\in\cD:\ell d\le n+\Delta\}.
 \label{eq:tpc36-column-upper-rank}
\end{align}
Then
\begin{equation}
 \one_{|\ell d-n|>\Delta}
 =1+H_{N_D}(k_-^\gamma(\ell),r(d))
    -H_{N_D}(k_+^\gamma(\ell),r(d)).
 \label{eq:tpc36-column-prefix}
\end{equation}

Choose \(M_D=2N_D+1\), and let
\(\widehat\kappa_{N_D}\) be as in
\eqref{eq:tpc36-interval-Fourier}.  Combining
\eqref{eq:tpc36-column-prefix} with
\eqref{eq:tpc36-prefix-Fourier} gives the pointwise identity
\begin{align}
 &\mathfrak m_{\rm act}((\ell,d),\gamma)
 \notag\\
 &\quad=\one_{\ell\ne\ell_\gamma}
 \one_{(d,e)\le G}
 \Bigg[1+\sum_{a\bmod M_D}
 \widehat\kappa_{N_D}(a)
 \bigl(
 e_{M_D}(a k_-^\gamma(\ell))
 -e_{M_D}(a k_+^\gamma(\ell))
 \bigr)e_{M_D}(-a r(d))\Bigg].
 \label{eq:tpc36-exact-column-Fourier}
\end{align}
Every factor outside the brackets is unary once \(\gamma\) is fixed.

\begin{theorem}[Exact fixed-column source--divisor separation]
\label{thm:tpc36-fixed-column-separation}
For every fixed opposite row \(\gamma\), there are coefficients and
bounded functions such that
\begin{equation}
 \mathfrak m_{\rm act}((\ell,d),\gamma)
 =\sum_{\rho\in\cR_\gamma}
 c_{\rho,\gamma}b_{\rho,\gamma}(\ell)
 f_{\rho,\gamma}(d),
 \label{eq:tpc36-column-separation}
\end{equation}
with
\begin{align}
 |\cR_\gamma|&\ll D,
 \label{eq:tpc36-column-layer-count}\\
 \sum_{\rho\in\cR_\gamma}|c_{\rho,\gamma}|
 \|b_{\rho,\gamma}\|_\infty
 \|f_{\rho,\gamma}\|_\infty
 &\ll1+\log(2D)=X^{o(1)}.
 \label{eq:tpc36-column-projective-cost}
\end{align}
The estimate is uniform in \(\gamma\).
\end{theorem}

\begin{proof}
Use the constant term and the \(M_D\) Fourier terms displayed in
\eqref{eq:tpc36-exact-column-Fourier}.  The source factor and the two
threshold phases belong to \(b\); the gcd cutoff and rank phase belong to
\(f\).  Their sup norms are bounded by \(2\) and \(1\), respectively.
The layer count is \(1+M_D=O(D)\), while
\eqref{eq:tpc36-interval-l1} proves the signed mass bound.  Importantly,
the normalized Fourier coefficients are summed in absolute value; the
raw layer count is not.
\end{proof}

\subsection{Including the physical row and orbit weights}

The coefficient \eqref{eq:tpc36-row-coefficients} is already a product of
a source weight and a \(d\)-weight, apart from a fixed-period factor.
First split the \(O_h(1)\) orbit cells \(j\bmod H_0\); finite Fourier
inversion in the remaining row residue variables has bounded total mass.
On each cell, Mellin inversion separates every smooth factor
\(W(\ell dj/X)\) into
\(\ell^{it}d^{it}j^{it}\) and has Schwartz total variation
\citep{WangTPC25}.  Multiplying these inherited representations by
\eqref{eq:tpc36-column-separation}, and enlarging \(\cR_\gamma\) by only
an \(O_h(1)\) factor, proves the physical form below.

\begin{theorem}[Physical product-ready column]
\label{thm:tpc36-physical-product-ready}
For each fixed \(\gamma\), the actual left-column multiplier including its
left row coefficient has the exact residue-cell representation
\begin{equation}
 \gamma_{(\ell,d)}^{(1)}
 \mathfrak A_{(\ell,d),\gamma}^{\rm act}(j)
 =\sum_{r\bmod H_0}\one_{j\equiv r\!\!\!\pmod{H_0}}
 \int_{\Omega_r}\sum_{\rho\in\cR_\gamma}
 c_{\rho,\omega,\gamma,r}
 b_{\rho,\omega,\gamma,r}(\ell)
 f_{\rho,\omega,\gamma,r}(d)
 \eta_{\rho,\omega,\gamma,r}(j/J)\,d\nu_r(\omega).
 \label{eq:tpc36-physical-product-ready}
\end{equation}
For every fixed \(B\ge0\), its normalized mass satisfies
\begin{equation}
 \sum_{r\bmod H_0}\int_{\Omega_r}\sum_{\rho\in\cR_\gamma}
 |c_{\rho,\omega,\gamma,r}|
 \|b_{\rho,\omega,\gamma,r}\|_\infty
 \|f_{\rho,\omega,\gamma,r}\|_\infty
 \|\eta_{\rho,\omega,\gamma,r}\|_{C^B}
 \,d|\nu_r|(\omega)
 \ll_B X^{o(1)},
 \label{eq:tpc36-product-ready-cost}
\end{equation}
after the fixed logarithmic source weight is included in the soft factor.
For fixed \(r\), the source function
\(b_{\rho,\omega,\gamma,r}(\ell)\) is independent of \(j\), \(d\), and
the complementary factor \(s\).
\end{theorem}

The statement is symmetric: fixing \(\alpha\) instead gives the analogous
right-column representation with the two row systems exchanged.

\begin{remark}[Coefficient normalization]
If the factor \(\log\ell\) is not absorbed into \(X^{o(1)}\), the right
side of \eqref{eq:tpc36-product-ready-cost} acquires the explicit harmless
factor \(O(\log L)\).  The theorem does not sum the physical row
coefficients in \(\ell^1\); it controls the multiplier representation
uniformly before the row sum.
\end{remark}

\subsection{Opening the ultra divisor}

Insert \eqref{eq:tpc36-physical-product-ready} into
\eqref{eq:tpc36-left-orbit-slice}, work on one of the fixed residue cells,
and open \eqref{eq:tpc36-ultra-increment}.  On each layer, all unary
support and smooth factors may be absorbed into \(b,f,g\), and
\eqref{eq:tpc36-reflection} gives
\begin{equation}
 \cZ_{\gamma,\rho,\omega}(j,s)
 =\sum_{\ell,d,u}
 b_{\gamma,\rho,\omega}(\ell)
 f_{\gamma,\rho,\omega}(d)
 g_{\gamma,\rho,\omega}(u)
 \one_{su=\ell jd+h}.
 \label{eq:tpc36-exact-integer-layer}
\end{equation}
This is the integer product-slope normal form.  Reducing
\eqref{eq:tpc36-exact-integer-layer} modulo \(q\) produces
\eqref{eq:tpc36-Rsa-interface}--\eqref{eq:tpc36-T-interface}; the source
weight remains independent of \(j,s\).

\begin{corollary}[The static half of the TPC-35 reassembly gate]
\label{cor:tpc36-static-half-gate}
The actual fixed-column mask, row coefficient, fixed-period factors, and
smooth orbit factors admit an exact \(X^{o(1)}\)-cost decomposition into
product-slope-ready layers.  What remains before the TPC-35 complete
finite-field spectrum can be applied to
\eqref{eq:tpc36-left-orbit-slice} is modular completion and coherent
alias reassembly, not static row-mask separation.
\end{corollary}

\subsection{Compatibility with the projected affine column}

TPC-33 alternatively opens the gcd projector after fixing
\(\gamma=(\ell_\gamma,e)\).  For \(v\mid e\), write \(d=vD'\).  Applying
\cref{thm:tpc36-fixed-column-separation} on each projected \(D'\)-range
gives at most
\begin{equation}
 \sum_{v\mid e}O\!\left(1+\frac Dv\right)
 \ll DX^{o(1)}
 \label{eq:tpc36-projected-layer-count}
\end{equation}
formal layers.  TPC-33 proves
\begin{equation}
 \sum_{v\mid e}|\lambda_G(v)|\le3^{\omega(e)},
 \label{eq:tpc36-fixed-column-gcd-mass}
\end{equation}
so their total signed Ferrers mass is
\begin{equation}
 \ll3^{\omega(e)}\log(2D)=X^{o(1)}.
 \label{eq:tpc36-projected-projective-cost}
\end{equation}
This version retains the affine determinant form of TPC-33, with product
slope \(\ell jv\).  Since \(v\ll D=o(J)\), multiplication by \(v\) is a
permutation modulo every prime \(q\asymp J\).  The direct version above is
shorter; the projected version confirms compatibility with the previous
M\"obius sign-fusion normal form.
